\documentclass[11pt]{amsart}
\usepackage[T1]{fontenc}
\usepackage{lmodern}
\usepackage{amsmath,amssymb,amsthm,mathtools}
\usepackage{booktabs}
\usepackage{xcolor}
\usepackage[colorlinks=true,linkcolor=blue!55!black,citecolor=blue!55!black,urlcolor=blue!55!black]{hyperref}

\newtheorem{theorem}{Theorem}[section]
\newtheorem{lemma}[theorem]{Lemma}
\newtheorem{proposition}[theorem]{Proposition}
\newtheorem{corollary}[theorem]{Corollary}
\theoremstyle{definition}
\newtheorem{definition}[theorem]{Definition}
\newtheorem{remark}[theorem]{Remark}

\newcommand{\CSF}{X}
\newcommand{\cycle}{C}

\title[Two transforms for unicyclic chromatic symmetric functions]
{Two structural transforms for chromatic symmetric functions of unicyclic graphs}
\author{Zijian Zeng}
\address{Institute of Computer Science and Digital Innovation, UCSI University,
Kuala Lumpur 56000, Malaysia}
\email{zijianzeng@foxmail.com}
\email{1002266693@ucsiuniversity.edu.my}
\date{August 30, 2026}
\hypersetup{
  pdftitle={Two structural transforms for chromatic symmetric functions of unicyclic graphs},
  pdfauthor={Zijian Zeng}
}

\begin{document}

\begin{abstract}
Let \(G\) be a connected unicyclic graph with cycle length \(c\).  We give
two linear transforms of the power-sum coefficients of its chromatic
symmetric function.  The first expresses a cycle-corrected degree polynomial;
the second expresses a cycle-corrected edge endpoint-degree polynomial.  If
two such graphs have the same chromatic symmetric function, their degree
polynomial difference is divisible by \(z(1-z)^c\), and by
\(z^2(1-z)^c\) when their leaf counts agree.  Combining this with known hook
coefficients proves the structural-data conjecture of Bingham, Johnston,
Lawson, Orellana, Pan, and Sato whenever the maximum degree is at most \(c\),
and throughout the infinite range \(c\geq4\), \(n\leq2c\).  Deterministic
canonical enumeration and independent exact modular evaluations also extend
the finite verification from order at most 16 through orders 17 and 18.  The
unrestricted conjecture remains open.
\end{abstract}

\maketitle

\section{Introduction}

Stanley's chromatic symmetric function \(\CSF_G\) refines the chromatic
polynomial by recording the color multiplicities of proper colorings
\cite{Stanley1995}.  Its ability to recover graph structure is particularly
subtle outside forests: nonisomorphic unicyclic graphs with equal chromatic
symmetric functions are known \cite{OrellanaScott2014}, including recent
examples of every cycle length at least four \cite{Bingham2026Girth}.

Bingham, Johnston, Lawson, Orellana, Pan, and Sato developed hook and leading
coefficient formulas for connected unicyclic graphs
\cite{BinghamEtAl2026}.  They conjectured that, for a triangle and odd order,
or for cycle length \(c\geq4\), equality of chromatic symmetric functions
forces equality of the degree sequence, the number of nontrivial rooted
trees attached to the cycle, and the number of internal edges.  The
version of record calls this Conjecture~63 and reports verification through
order 16.  (It is Conjecture~9.1 in the first arXiv version.)

Our main device is a binomial transform of the power-sum coefficients weighted
by the number of singleton parts.  For a tree, the transform is the ordinary
degree polynomial.  For a unicyclic graph it separates vertices on the cycle
from vertices off the cycle by the universal kernel
\(1-(1-z)^{c-1}\).  This yields a high-multiplicity zero at \(z=1\) for
the difference of degree polynomials of equal-CSF graphs.  A companion
transform weighted by the number of parts of size two gives an analogous
formula for the edge endpoint-degree polynomial.

The resulting consequences are genuine partial results rather than a proof
of Conjecture~63.  The remaining obstruction has two aspects: degree
multiplicities above the cycle length can support a nonzero divisible
difference, and the known hook formulas leave an exceptional alternative
with one nontrivial rooted tree on one graph and \(c\) on the other.

\section{Definitions and published reductions}

All graphs are finite, undirected, and simple.  A connected graph is
\emph{unicyclic} if it has exactly one cycle.  Let \(G=(V,E)\) be such a
graph, let \(n=|V|=|E|\), and let \(\cycle\) be its unique cycle, of length
\(c\).  Deleting the cycle edges leaves a rooted tree \(T_v\) at each
\(v\in V(\cycle)\).  Write
\[
r(G)=\#\{v\in V(\cycle): |V(T_v)|>1\}
    =\#\{v\in V(\cycle):\deg_G(v)\geq3\}.
\]
An edge is \emph{internal} if both endpoints have degree at least two; let
\(k(G)\) be the number of internal edges.  If \(L(G)\) is the number of
leaves, then every non-internal edge is incident with a unique leaf, and
therefore
\begin{equation}\label{eq:k-leaves}
k(G)=n-L(G).
\end{equation}

Write the power-sum expansion as
\[
\CSF_G=\sum_{\lambda\vdash n}a_\lambda(G)p_\lambda.
\]
For a partition \(\lambda\), let \(m_j(\lambda)\) be the multiplicity of
the part \(j\), and let \(\ell(\lambda)\) be its length.  We use Stanley's
all-subsets formula
\begin{equation}\label{eq:stanley}
\CSF_G=\sum_{F\subseteq E}(-1)^{|F|}p_{\lambda(V,F)},
\end{equation}
where \(\lambda(V,F)\) records the component orders of \((V,F)\)
\cite[Theorem~2.5]{Stanley1995}.

We use two results from \cite{BinghamEtAl2026}.  First, equal chromatic
symmetric functions of connected unicyclic graphs have the same cycle length.
Second, their hook coefficients force either
\begin{equation}\label{eq:hookordinary}
(r(G),k(G))=(r(H),k(H)),
\end{equation}
or, after interchanging the graphs,
\begin{equation}\label{eq:hookexceptional}
r(G)=1,\qquad r(H)=c,\qquad k(G)=k(H)+1.
\end{equation}
This is Proposition~16 of that paper.

\section{The degree transform}

Define
\[
s_b(G)=(-1)^{n-b}\sum_{\lambda\vdash n}a_\lambda(G)
m_1(\lambda)\binom{n-\ell(\lambda)}{n-b},
\qquad
\mathcal S_G(z)=\sum_{b=0}^{n}s_b(G)z^b,
\]
where an out-of-range binomial coefficient is zero.  Separate the degree
polynomial into its cycle and off-cycle parts:
\begin{align*}
A_G(z)&=\sum_{v\in V(\cycle)}z^{\deg(v)},&
O_G(z)&=\sum_{v\notin V(\cycle)}z^{\deg(v)},\\
D_G(z)&=A_G(z)+O_G(z).&&
\end{align*}

\begin{theorem}[Degree transform]\label{thm:degree-transform}
For every connected unicyclic graph \(G\),
\begin{equation}\label{eq:degree-transform}
\mathcal S_G(z)=A_G(z)+\bigl(1-(1-z)^{c-1}\bigr)O_G(z).
\end{equation}
\end{theorem}

\begin{proof}
Substitute \eqref{eq:stanley} into the definition of \(s_b\).
The factor \(m_1(\lambda(V,F))\) permits a double count over isolated
vertices.  Fix such a vertex \(v\), put \(q=n-b\), and let
\(\kappa(F)\) be the number of components.  Its contribution is
\begin{equation}\label{eq:fixed-vertex}
\sum_{F\subseteq E(G-v)}(-1)^{|F|+q}
\binom{n-\kappa(F)}q.
\end{equation}

We first dispose of \(q=0\), or \(b=n\).  The fixed-vertex contribution is
\(\sum_{F\subseteq E(G-v)}(-1)^{|F|}=0\), because \(G-v\) has at least one
edge: deleting a cycle vertex leaves an edge of the broken cycle, while
deleting an off-cycle vertex leaves the whole cycle.  The right side of
\eqref{eq:degree-transform} also has degree at most \(n-1\).  Hence assume
\(q\geq1\) below.

Suppose first that \(v\in V(\cycle)\).  Every \(F\subseteq E(G-v)\) is a
forest, so \(n-\kappa(F)=|F|\).  Since
\(|E(G-v)|=n-\deg(v)\), binomial inversion gives
\[
\sum_{F\subseteq E(G-v)}(-1)^{|F|+q}\binom{|F|}q
=\boldsymbol 1_{\{\deg(v)=b\}}.
\]
Thus \(v\) contributes \(z^{\deg(v)}\).

Now suppose that \(v\notin V(\cycle)\), put \(d=\deg(v)\), and let
\(h=n-d-c\) be the number of non-cycle edges in \(G-v\).  The same forest
calculation supplies the base term, but is wrong precisely for edge sets
containing all of \(E(\cycle)\).  For such a set,
\(n-\kappa(F)=|F|-1\).  The correction to
\eqref{eq:fixed-vertex} is
\begin{align*}
&-(-1)^{c+q}\sum_{j=0}^{h}(-1)^j\binom hj
  \binom{c+j-1}{q-1}\\
&\qquad=(-1)^{b+d+1}\binom{c-1}{d+c-b-1},
\end{align*}
where the last equality uses
\[
\sum_{j=0}^{h}(-1)^j\binom hj\binom{u+j}{q}
=(-1)^h\binom u{q-h}.
\]
The \(j=0\) correction cancels the forest base term.  The remaining
polynomial contribution of \(v\) is
\[
z^d\sum_{j=1}^{c-1}(-1)^{j+1}\binom{c-1}{j}z^j
=z^d\bigl(1-(1-z)^{c-1}\bigr).
\]
Summing over all vertices proves the identity.
\end{proof}

\begin{corollary}\label{cor:degree-divisibility}
If \(G\) and \(H\) are connected unicyclic graphs with
\(\CSF_G=\CSF_H\), and their common cycle length is \(c\), then
\begin{equation}\label{eq:degree-divisibility}
D_G(z)-D_H(z)=(1-z)^{c-1}\bigl(O_G(z)-O_H(z)\bigr)
\end{equation}
and
\begin{equation}\label{eq:basic-divisibility}
z(1-z)^c\mid D_G(z)-D_H(z).
\end{equation}
If \(L(G)=L(H)\), then
\begin{equation}\label{eq:leaf-divisibility}
z^2(1-z)^c\mid D_G(z)-D_H(z).
\end{equation}
\end{corollary}

\begin{proof}
Equality in \eqref{eq:degree-transform} yields
\eqref{eq:degree-divisibility}.  The polynomial \(O_G-O_H\) has zero
constant term and vanishes at \(z=1\), since both graphs have \(n-c\)
off-cycle vertices.  This proves \eqref{eq:basic-divisibility}.  If the leaf
counts agree, then its coefficient of \(z\) is also zero, which proves
\eqref{eq:leaf-divisibility}.
\end{proof}

The coefficient of \(z^2\) in Theorem~\ref{thm:degree-transform} recovers
the invariant \((c-1)L+I_\cycle\) of Orellana and Scott, where
\(I_\cycle\) counts degree-two cycle vertices
\cite[Proposition~4.1]{OrellanaScott2014}.  The full polynomial identity
retains every degree.

\section{Structural consequences}

\begin{theorem}\label{thm:structural-range}
Let \(G,H\) be connected unicyclic graphs with equal chromatic symmetric
functions and common cycle length \(c\).
\begin{enumerate}
\item If \(\max\{\Delta(G),\Delta(H)\}\leq c\), then \(G,H\) have the same
degree sequence, the same number of nontrivial rooted trees, and the same
number of internal edges.
\item Conjecture~63 holds for every \(c\geq4\) and \(n\leq2c\).
\end{enumerate}
\end{theorem}

\begin{proof}
For the first assertion, a nonzero polynomial divisible by
\(z(1-z)^c\) has degree at least \(c+1\).  Thus
\eqref{eq:basic-divisibility} forces \(D_G=D_H\).  Equation
\eqref{eq:k-leaves} gives equal internal-edge counts, and the published hook
formula then gives equal rooted-tree counts (with the cycle case detected
separately by its hook coefficient).

For the second assertion, first suppose \(n\leq2c-1\).  The exceptional
alternative \eqref{eq:hookexceptional} is impossible: \(r(H)=c\) requires
at least one off-cycle vertex at each cycle vertex and hence \(n\geq2c\).
Thus \eqref{eq:hookordinary} holds, so the leaf counts agree by
\eqref{eq:k-leaves}.  Every connected \(c\)-unicyclic graph satisfies
\begin{equation}\label{eq:max-degree}
\Delta(G)\leq n-c+2\leq c+1.
\end{equation}
Indeed, a vertex has at most all \(n-c\) off-cycle incidences together with
two cycle incidences.  Since a nonzero polynomial divisible by
\(z^2(1-z)^c\) has degree at least \(c+2\),
\eqref{eq:leaf-divisibility} gives \(D_G=D_H\).

It remains to treat \(n=2c\).  In the ordinary hook alternative, the leaf
counts again agree, and \eqref{eq:max-degree} gives degree at most \(c+2\).
Therefore
\[
D_G-D_H=a z^2(1-z)^c
\]
for some integer \(a\).  If \(c\geq5\) and \(a\ne0\), the coefficient
\(\ell^1\)-norm is at least \(2^c>4c=2n\), but the difference of two
\(n\)-vertex degree polynomials has norm at most \(2n\), a contradiction.
For \(c=4\), only \(|a|=1\) survives the norm bound.  Equality in that bound
would force, on one side, the degree multiset
\((6,4,4,4,4,4,4,2)\), whose degree sum is 32 rather than \(2n=16\).

Finally assume the exceptional alternative.  The graph \(H\) with \(r(H)=c\)
has exactly one off-cycle vertex at each cycle vertex, so it is the
\(c\)-sun.  The graph \(G\) has \(c-1\) leaves and a unique off-cycle
internal vertex.  That internal vertex needs a neighbor other than the cycle
root, so the root is incident with at most \(c-2\) leaves and
\(\Delta(G)\leq c+1\).  Divisibility, the degree bound, and the coefficient
of \(z\) force
\[
D_G-D_H=-z(1-z)^c.
\]
Comparison at \(z^2\) forces the off-cycle internal vertex to have degree
two and the unique nontrivial cycle root to have degree \(c+1\).  The actual
difference is then
\[
-z+cz^2-cz^3+z^{c+1}.
\]
Equality with \(-z(1-z)^c\) would imply \(\binom c2=c\), hence \(c=3\),
contrary to \(c\geq4\).  This completes the proof.
\end{proof}

\begin{remark}
The excluded value \(c=3,n=6\) is not an artifact: the last two displayed
polynomials then agree, matching the known even-order triangle phenomenon.
\end{remark}

\section{The edge endpoint-degree transform}

Define
\[
t_b(G)=(-1)^{n-b}\sum_{\lambda\vdash n}a_\lambda(G)
m_2(\lambda)\binom{n-\ell(\lambda)}{n-b},
\qquad \mathcal T_G(z)=\sum_{b=0}^{n}t_b(G)z^b.
\]
Let
\[
P_G(z)=\sum_{uv\in E(G)}z^{\deg(u)+\deg(v)-2},
\]
and let \(Q_G(z)\) be the same sum restricted to edges with both endpoints
off \(\cycle\).

\begin{theorem}[Edge transform]\label{thm:edge-transform}
For every connected unicyclic graph \(G\),
\begin{equation}\label{eq:edge-transform}
\mathcal T_G(z)=(1-z)P_G(z)-(1-z)^{c}Q_G(z).
\end{equation}
Consequently, equal-CSF connected unicyclic graphs of cycle length \(c\)
satisfy
\begin{equation}\label{eq:edge-divisibility}
z(1-z)^{c-1}\mid P_G(z)-P_H(z).
\end{equation}
\end{theorem}

\begin{proof}
In \eqref{eq:stanley}, \(m_2(\lambda(V,F))\) counts two-vertex components.
Fix such a component on an edge \(e=uv\), put
\(d_e=\deg(u)+\deg(v)-2\), and select an arbitrary edge set \(R\) in
\(G-\{u,v\}\).  All other edges incident with \(u\) or \(v\) must be
absent.  There are \(M=n-d_e-1\) available edges.

If \(R\) is a forest, its contribution for \(q=n-b\) is
\((-1)^{|R|+1+q}\binom{|R|+1}{q}\).  The identity
\[
\sum_{r=0}^{M}(-1)^r\binom Mr(1+x)^{r+1}=(1+x)(-x)^M
\]
shows that the fixed edge contributes \(z^{d_e}(1-z)\).
If \(u\) or \(v\) is on \(\cycle\), all relevant \(R\) are forests.  If
both are off \(\cycle\), the edge sets containing the whole cycle require a
correction.  With \(h=n-d_e-c-1\), finite differences give
\[
\sum_{j=0}^{h}(-1)^{c+j+q}\binom hj\binom{c+j}{q-1}
=(-1)^{b+d_e+1}\binom c{d_e+c-b},
\]
whose generating polynomial is \(-z^{d_e}(1-z)^c\).  Summing over edges
proves \eqref{eq:edge-transform}.

The quotient \(\mathcal T_G/(1-z)\) is determined by \(\CSF_G\).  Thus
\[
P_G-P_H=(1-z)^{c-1}(Q_G-Q_H).
\]
An edge entirely off the cycle cannot have two degree-one endpoints in a
connected unicyclic graph, so \(Q_G(0)=Q_H(0)=0\).  This proves
\eqref{eq:edge-divisibility}.
\end{proof}

For trees, the generalized degree polynomial determined by the chromatic
symmetric function recovers richer edge-degree information
\cite{AlistePrietoEtAl2024}.  Theorem~\ref{thm:edge-transform} is instead a
direct unicyclic formula with an explicit cycle correction.

\section{Finite verification at orders 17 and 18}

We used \texttt{geng} from nauty~2.7r3 to generate one representative of each
isomorphism class of connected simple graphs on \(n\) vertices and exactly
\(n\) edges \cite{McKayPiperno2014}.  Such a graph is unicyclic.  The totals
agree with the independent enumeration A001429 \cite{OEISA001429}.

For a graph decomposed into rooted trees at its cycle, choose fixed values
\(x_1,\ldots,x_q\) and put \(B=J-I\).  We evaluated
\begin{equation}\label{eq:trace-eval}
\CSF_G(x_1,\ldots,x_q)=
\operatorname{tr}(D_0B D_1B\cdots D_{c-1}B),
\end{equation}
where the diagonal matrix \(D_i\) contains the rooted-tree coloring weights
at the \(i\)-th cycle vertex.  All arithmetic was exact modulo fixed primes.
Equal symmetric functions necessarily give equal values in every such
specialization, so unequal values are deterministic certificates of
inequality.  Equal values were treated only as collisions, never as proofs of
CSF equality.

Two implementations were used.  The discovery implementation has a custom
graph6 parser, leaf-pruned cycle extraction, and a start-color recurrence.  A
separate verifier uses NetworkX parsing and cycle bases
\cite{HagbergEtAl2008}, recursive rooted-tree messages, a matrix-trace
aggregation, and disjoint moduli and evaluation points.  A prospective pair
with different structural data surviving all evaluations would be passed to
two full integer power-sum implementations: direct spanning-subset
enumeration and vertex-weighted deletion--contraction.

\begin{proposition}\label{prop:finite}
No two connected unicyclic graphs of order 17 have the same chromatic
symmetric function.  At order 18, Conjecture~63 holds for every graph in its
stated domain.
\end{proposition}

\begin{proof}[Computational certificate]
The coverage and collision counts are as follows.
\begin{center}
\small
\begin{tabular}{@{}rrrrr@{}}
\toprule
\(n\) & generated & included & collisions & survivors\\
\midrule
17 & 880,840 & 880,840 & 0 & 0\\
18 & 2,497,405 & 1,671,621 & 1 & 0\\
\bottomrule
\end{tabular}
\end{center}
At order 17 every primary signature is distinct, proving the first statement.
At order 18 the sole primary collision also matched at the independent
refinement value, but the two records had the same degree sequence,
rooted-tree count, and internal-edge count.  Thus no equal-CSF pair can have
different structural data.  We make no claim that this same-data evaluation
collision represents equal chromatic symmetric functions.

At order 18 the conjecture domain omits the even-order triangle class.  The
cycle-length subtotals for the 1,671,621 included graphs sum over
\(c=4,\ldots,18\), and the generator's complete total is 2,497,405.  These
checks, together with the canonical-generation guarantee, establish the
claimed coverage.
\end{proof}

Combining Proposition~\ref{prop:finite} with the verification through order
16 reported in \cite{BinghamEtAl2026} verifies Conjecture~63 through order 18.

\section{Limitations and further questions}

For equal-CSF graphs, Corollary~\ref{cor:degree-divisibility} writes
\(D_G-D_H\) as a high-order finite difference.  When degrees exceed \(c\),
the quotient need not vanish on purely algebraic grounds.  Even when the
published hook coefficients give equal leaf counts, the possible factor
\(z^2(1-z)^{c}R(z)\) has room for a nonconstant integer polynomial once
\(n>2c\).  The edge transform adds a different necessary condition but does
not yet force this quotient to vanish.

The other unresolved branch is \eqref{eq:hookexceptional}.  Feasibility alone
forces \(n\geq2c\), and Theorem~\ref{thm:structural-range} eliminates its
endpoint for \(c\geq4\), but no contradiction is known here for general
\(n>2c\).  A transform based on components of size three or a more local use
of the leading coefficient may provide the next constraint.  Exhaustive
search beyond order 18 remains evidence only, not a substitute for such an
argument.

\section*{Data, code, and computational disclosure}

The accompanying archive contains the literal source statement, the formal
statement, proofs, search records, exact test programs, the two independent
stream evaluators, fixed finite-result files, and commands for reproduction.
The principal fast tests compare three exact CSF computations and check both
transform identities.  A full small-order regression covers all 1,040
connected unicyclic graphs of orders 3 through 10.

Computer algebra and custom exact programs were used during exploratory
algebra and implementation.  All proofs were checked from definitions, and
all finite claims were reproduced by deterministic exact computation.  The
author takes responsibility for the results and exposition.

\bibliographystyle{amsplain}
\bibliography{references}

\end{document}
